\chapter{Clase 14: Sistemas Jerárquicos (parte 2)}

Esta clase es completamente práctica y se desarrolla íntegramente en el cuaderno de Jupyter \texttt{CuadernoClase\_SistemasJerarquicos\_Ejercicio(1).ipynb}. El objetivo es aplicar todo lo visto en la Clase 13: construir explícitamente un sistema jerárquico de tres cuerpos usando coordenadas de Jacobi y simular su evolución dinámica con la librería \texttt{pymcel}.

\section{Definición del sistema jerárquico de tres cuerpos}
Llamamos al sistema completo \textbf{Zul}. Sus componentes son:
\begin{itemize}
    \item \textbf{Zul Aa}: estrella primaria (\(m = 1.0\))
    \item \textbf{Zul Ab}: planeta (\(m = 0.01\))
    \item \textbf{Zul B}: estrella secundaria (\(m = 5.0\))
\end{itemize}

El árbol jerárquico es:

\begin{itemize}[leftmargin=2em,label={}]
\item \textbf{Zul} (sistema total)
  \begin{itemize}[leftmargin=2em,label={}]
  \item \textbf{A} (subsistema interno: Aa + Ab)
    \begin{itemize}[leftmargin=2em,label={}]
    \item \textbf{Aa} (estrella primaria, \(m_{\text{Aa}} = 1.0\))
    \item \textbf{Ab} (planeta, \(m_{\text{Ab}} = 0.01\))
    \end{itemize}
  \item \textbf{B} (estrella secundaria, \(m_{\text{B}} = 5.0\))
  \end{itemize}
\end{itemize}

\begin{importante}
Este es un sistema jerárquico típico: un planeta orbita una estrella (subsistema A) y ese subsistema orbita a su vez una estrella más masiva (B). Las masas y distancias están elegidas para que la jerarquía sea estable.
\end{importante}

\section{Condiciones iniciales en coordenadas de Jacobi}
Usamos unidades canónicas: \(G = 1\).

\subsection{Subsistema interno A (Aa + Ab)}
\[
m_{\text{A}} = m_{\text{Aa}} + m_{\text{Ab}} = 1.01
\]
\[
\vec{r}_{\text{A}} = \begin{pmatrix} 0.1 \\ 0 \\ 0 \end{pmatrix}, \quad
\vec{v}_{\text{A}} = \begin{pmatrix} 0 \\ 4.0 \\ 0 \end{pmatrix}
\]

\subsection{Vector relativo entre A y B}
\[
\vec{r}_{\text{AB}} = \begin{pmatrix} 2.0 \\ 0 \\ 0 \end{pmatrix}, \quad
\vec{v}_{\text{AB}} = \begin{pmatrix} 0 \\ 2.0 \\ 0 \end{pmatrix}
\]

\subsection{Centro de masa total del sistema}
\[
\vec{R}_{\text{CM}} = \begin{pmatrix} 0 \\ 0 \\ 0 \end{pmatrix}, \quad
\vec{V}_{\text{CM}} = \begin{pmatrix} 0.01 \\ 0 \\ 0 \end{pmatrix}
\]

\section{Descenso por el árbol jerárquico (cálculo paso a paso)}

\subsection{Nivel 1: Posiciones y velocidades del CM de A y de B}
\[
\vec{R}_{\text{CM,A}} = \vec{R}_{\text{CM}} + \frac{m_{\text{B}}}{m_{\text{A}} + m_{\text{B}}} \vec{r}_{\text{AB}} = \begin{pmatrix} 1.66389351 \\ 0 \\ 0 \end{pmatrix}
\]
\[
\vec{V}_{\text{CM,A}} = \vec{V}_{\text{CM}} + \frac{m_{\text{B}}}{m_{\text{A}} + m_{\text{B}}} \vec{v}_{\text{AB}} = \begin{pmatrix} 0.01 \\ 1.66389351 \\ 0 \end{pmatrix}
\]
\[
\vec{r}_{\text{B}} = \vec{R}_{\text{CM}} - \frac{m_{\text{A}}}{m_{\text{A}} + m_{\text{B}}} \vec{r}_{\text{AB}} = \begin{pmatrix} -0.33610649 \\ 0 \\ 0 \end{pmatrix}
\]
\[
\vec{v}_{\text{B}} = \vec{V}_{\text{CM}} - \frac{m_{\text{A}}}{m_{\text{A}} + m_{\text{B}}} \vec{v}_{\text{AB}} = \begin{pmatrix} 0.01 \\ -0.33610649 \\ 0 \end{pmatrix}
\]

\subsection{Nivel 2: Posiciones y velocidades de Aa y Ab dentro de A}
\[
\vec{r}_{\text{Aa}} = \vec{R}_{\text{CM,A}} + \frac{m_{\text{Ab}}}{m_{\text{A}}} \vec{r}_{\text{A}} = \begin{pmatrix} 1.66488361 \\ 0 \\ 0 \end{pmatrix}
\]
\[
\vec{v}_{\text{Aa}} = \vec{V}_{\text{CM,A}} + \frac{m_{\text{Ab}}}{m_{\text{A}}} \vec{v}_{\text{A}} = \begin{pmatrix} 0.01 \\ 1.70349747 \\ 0 \end{pmatrix}
\]
\[
\vec{r}_{\text{Ab}} = \vec{R}_{\text{CM,A}} - \frac{m_{\text{Aa}}}{m_{\text{A}}} \vec{r}_{\text{A}} = \begin{pmatrix} 1.56488361 \\ 0 \\ 0 \end{pmatrix}
\]
\[
\vec{v}_{\text{Ab}} = \vec{V}_{\text{CM,A}} - \frac{m_{\text{Aa}}}{m_{\text{A}}} \vec{v}_{\text{A}} = \begin{pmatrix} 0.01 \\ -2.29650253 \\ 0 \end{pmatrix}
\]

\begin{importante}
Estas fórmulas son las coordenadas de Jacobi aplicadas recursivamente. Cada nivel usa el centro de masa del subsistema anterior como nuevo origen.
\end{importante}

\section{Construcción del sistema para simulación}
El sistema completo se define como una lista de diccionarios:
\begin{lstlisting}
sistema = [
    dict(m = 1.0,   r = r_Zul_Aa, v = v_Zul_Aa),
    dict(m = 0.01,  r = r_Zul_Ab, v = v_Zul_Ab),
    dict(m = 5.0,   r = r_Zul_B,  v = v_Zul_B)
]
\end{lstlisting}

Se simula con la función \texttt{ncuerpos\_solución} de \texttt{pymcel} en el intervalo temporal:
\[
t \in [0, 50]\qquad \text{con 1000 puntos}
\]

\section{Resultados de la simulación y visualización}
La simulación muestra claramente el comportamiento jerárquico esperado:
\begin{itemize}
    \item \textbf{Órbita interna rápida}: el planeta Ab describe una órbita casi circular alrededor de Aa (radio \(\approx 0.1\)).
    \item \textbf{Órbita intermedia}: el subsistema A (Aa + Ab) describe una órbita alrededor del centro de masa total.
    \item \textbf{Órbita externa}: la estrella B orbita alrededor del centro de masa del sistema completo con un radio mayor (\(\approx 0.336\)).
\end{itemize}

Se generan gráficas de:
\begin{itemize}
    \item Trayectorias en el plano \(xy\) (todas las partículas).
    \item Trayectorias en el marco del centro de masa.
    \item Animación con \texttt{celluloid} que muestra el movimiento jerárquico completo.
\end{itemize}

\begin{importante}
La simulación numérica confirma que el sistema permanece estable durante el tiempo integrado. Esto valida la aproximación jerárquica y la correcta inicialización con coordenadas de Jacobi.
\end{importante}

\section{Generalización y consejos prácticos}
\begin{enumerate}
    \item Para cualquier sistema jerárquico de \(N\) cuerpos se sigue el mismo procedimiento recursivo: se desciende por el árbol calculando centros de masa y vectores relativos en cada nivel.
    \item Siempre se debe fijar primero el centro de masa y velocidad total del sistema (\(\vec{R}_{\text{CM}}\), \(\vec{V}_{\text{CM}}\)).
    \item La librería \texttt{pymcel} facilita enormemente la simulación una vez que se tienen las condiciones iniciales en el formato de lista de diccionarios.
    \item En sistemas reales (ej. sistemas estelares múltiples, exoplanetas circumbinarios) se usan exactamente estas mismas técnicas para generar condiciones iniciales estables.
\end{enumerate}

\section{Conclusiones de las Clases 13 y 14}
Las dos clases juntas muestran el camino completo:
\begin{itemize}
    \item Teoría (Clase 13): motivación, autoorganización y coordenadas de Jacobi.
    \item Práctica (Clase 14): construcción explícita de un sistema jerárquico y simulación numérica.
\end{itemize}

Esta metodología es la base para estudiar sistemas de \(N\) cuerpos más complejos en las próximas unidades.



